%==============================================================================
% IaCSecBench -- manuscript prepared for Empirical Software Engineering.
%
% CLASS. EMSE requires Springer's svjour3 class. Springer distributes it only in
% its own template archive: it is not on CTAN and not in any TeX distribution's
% package set, so a build cannot fetch it. Download the template from the
% journal's submission-guidelines page, put svjour3.cls (with svglov3.clo and
% spbasic.bst) beside this file, and the conditional below switches to it with no
% further edit to the manuscript.
%
% Without those files this builds under `article' at a text measure deliberately
% narrower than Springer's smallextended, so a table that fits here fits there.
% Every table is set with tabularx against \linewidth for the same reason.
%
% CITATIONS. EMSE uses author-year. \citep and \citet are used throughout and
% bare \cite nowhere, because bare \cite means different things to natbib in
% author-year and numeric modes.
%==============================================================================

\newif\ifhassvjour
\IfFileExists{svjour3.cls}{\hassvjourtrue}{\hassvjourfalse}

\ifhassvjour
  \documentclass[smallextended]{svjour3}
\else
  \documentclass[11pt,a4paper]{article}
\fi

\usepackage{amsmath,amssymb}
\usepackage{graphicx}
\usepackage{booktabs}
\usepackage{tabularx}
\usepackage{array}

% natbib may already be loaded by the journal class; loading it twice with
% different options is an error, so ask first.
\makeatletter
\@ifpackageloaded{natbib}{}{\usepackage{natbib}}
\makeatother

\ifhassvjour\else
  % Fallback-only scaffolding. The measure is narrower than Springer's, so this
  % is the conservative direction: nothing that fits here overflows there.
  \usepackage[a4paper,textwidth=117mm,textheight=205mm]{geometry}
  \newcommand{\keywords}[1]{\par\smallskip\noindent\textbf{Keywords}\quad #1}
  \newcommand{\institute}[1]{}
  \newcommand{\journalname}[1]{}
  \providecommand{\and}{\unskip\ $\cdot$\ \ignorespaces}
  \newcommand{\titlerunning}[1]{}
  \newcommand{\authorrunning}[1]{}
\fi

% Long file paths set in typewriter are single unbreakable tokens, and at a
% one-column measure they overflow the line rather than wrapping. \path (from the
% url package, which hyperref loads) breaks after `/' without inserting a hyphen,
% which in a path would be read as part of the path, and needs no escaping of `_',
% so a path cannot acquire a stray backslash in transcription.
\usepackage[hidelinks,breaklinks]{hyperref}   % load last

% EMSE asks for DOIs "as full DOI links in your reference list (e.g.
% 'https://doi.org/abc')". spbasic prints "DOI 10.xxx" and defines \doi only if
% it is undefined, so defining it here wins. The argument is read with `_' as an
% ordinary character because one DOI in refs.bib contains an underscore.
\makeatletter
\providecommand{\doi}{}
\renewcommand{\doi}{\begingroup\catcode`\_=12\relax\iacsb@doi}
\def\iacsb@doi#1{\endgroup\href{https://doi.org/#1}{\nolinkurl{https://doi.org/#1}}}
\makeatother

\providecommand{\orcidID}[1]{\unskip~(\href{https://orcid.org/#1}{\footnotesize ORCID: #1})}

% Author-year citation groups are long, and two or three in one parenthesis do not
% fit a narrow measure at any breakpoint TeX would otherwise choose.
% \emergencystretch is LaTeX's sanctioned last resort for exactly this: it widens
% interword space only on lines that would otherwise be overfull.
\emergencystretch=2em

% Tables are generated artefacts; this is where they come from. Numeric columns
% in every emitter are `r', never `c': each column is written to a fixed number
% of decimal places, so right alignment puts the decimal points of 3.85 and 76.92
% in the same place and centring does not. siunitx is deliberately absent -- the
% fixed decimal places make it unnecessary, and its option names differ between
% versions 2 and 3.
\graphicspath{{figures/}{../figures/}}
\newcommand{\resulttable}[1]{%
  \IfFileExists{../results/tables/#1.tex}%
    {\input{../results/tables/#1.tex}}%
    {\begin{table}[!t]\centering\small
       \fbox{\parbox{0.9\linewidth}{\centering
         \texttt{\detokenize{results/tables/#1.tex}} not found.\\[2pt]
         Run \path{experiments/run_baselines.sh} to generate it.}}%
     \end{table}}}

% A figure that does not yet exist prints a visible placeholder rather than
% breaking the build, so an incomplete figure set cannot pass silently as
% finished work. \detokenize re-catcodes `_' so call sites need not escape it.
%
% Both drawings are generated: experiments/generate_figures.py writes their
% Mermaid sources from results/evaluation.json and results/run_manifest.json, so
% a figure cannot assert a corpus size or tool version that disagrees with a
% table. Re-render with the mmdc command in paper/figures/README.md after any
% re-measurement; CI fails if the committed sources go stale.
\newcommand{\paperfigure}[3]{%
  \begin{figure}[!t]\centering
    \IfFileExists{#1}{\includegraphics[width=\linewidth]{#1}}%
      {\fbox{\parbox{0.9\linewidth}{\centering\vspace{2em}
        \texttt{\detokenize{#1}}\\[4pt] figure not yet produced\vspace{2em}}}}
    \caption{#2}\label{#3}
  \end{figure}}

\journalname{Empirical Software Engineering}

\begin{document}

\title{IaCSecBench: A Finding-Normalization Methodology and Reproducible
Harness for Evaluating Infrastructure-as-Code Security Validation}
% svjour3 warns that the full title is too long for the running head.
\titlerunning{IaCSecBench: Finding Normalization for IaC Security Validation}

\author{Manideep Chittineni \orcidID{0009-0003-9709-5842}}
\authorrunning{M. Chittineni}

\institute{M. Chittineni \at
  Independent Researcher, London, United Kingdom \\
  \email{manideep.chittineni@hotmail.com}}

\date{Received: date / Accepted: date}

\maketitle

%------------------------------------------------------------------------------
\begin{abstract}
\textbf{Context.} Infrastructure as Code (IaC) dominates cloud provisioning, yet pre-deployment security validation remains fragmented across static scanners, policy engines, and test frameworks. Incompatible abstraction levels, bespoke schemas, and developer-authored benchmarks impede rigorous comparison.
\textbf{Objective.} We investigate how far the matching criterion (what counts as a detection) determines a comparison's conclusions.
\textbf{Method.} IaCSecBench is an open evaluation harness that normalizes heterogeneous scanner outputs onto a canonical control taxonomy, enforces mechanical corpus admissibility, and scores under three matching criteria. Over 56 admissible cases (30 vulnerable, 26 compliant), we evaluate three third-party scanners, plan-level policy evaluation, and a repository-edge layer, reporting exact Clopper--Pearson intervals and exact McNemar tests under Holm--Bonferroni correction.
\textbf{Results.} Relaxing the criterion from control attribution to any finding shifts apparent recall by 6.7 to 10.0 percentage points for source- and plan-level tools, under half the 23.3-point spread between them, and by 90.0 points for repository-edge scanning, promoting it from last to second. No difference among the four survives correction; under it, three of their six pairs could not have survived under any outcome, and the rest only with a unanimous split. A third-party scanner, not our pipeline, attains the highest point-estimate recall.
\textbf{Conclusion.} A comparison that does not state its matching criterion cannot be interpreted, since that choice alone moved one layer from last to second. We release our open replication package; every results table is a generated artefact.
\end{abstract}

\keywords{Infrastructure as Code \and Security validation \and Benchmarking
methodology \and Static analysis \and Policy as code \and Empirical evaluation}

%------------------------------------------------------------------------------
\section{Introduction}
\label{sec:intro}

Cloud platforms are increasingly provisioned through Infrastructure as Code
(IaC), in which infrastructure topology is expressed declaratively in
version-controlled artefacts \citep{morris2020infrastructure,
guerriero2019adoption}, atop service models whose elasticity makes such
automation a practical necessity \citep{armbrust2010view}. The property that
makes IaC attractive, deterministic and repeatable application of a declared
state, also amplifies its failure modes: a single insecure declaration, such as
an unrestricted ingress rule or a disabled encryption setting, propagates
automatically to every environment that consumes the module. Empirical studies of
configuration repositories confirm that such defects are common, recurring as
hard-coded secrets, excessive privileges, and insecure defaults
\citep{rahman2019sins, rahman2021replication, verdet2023exploring,
verdet2025adoption}. \citet{rahman2019mapping} map the resulting research area
and identify tooling and empirical evaluation as two of its four constituent
topics.

Because remediation after provisioning implies a window of exposure,
contemporary practice shifts validation earlier, into continuous integration
\citep{bass2015devops, humble2010continuous}. Doing so raises an evaluation
problem the literature has not resolved. Candidate gates operate at incomparable
abstraction levels: lexical scanning of HashiCorp Configuration Language (HCL) source, abstract-syntax-tree
analysis, and policy evaluation over resolved plan state. Each reports findings
in a bespoke schema, with rule identifiers that carry no shared semantics.
Comparative claims are consequently difficult to reproduce, and when a benchmark
is authored by the developer of one of the tools under comparison, agreement
between that tool and the ground truth is close to guaranteed by construction.

This paper addresses the evaluation problem; the pipeline it evaluates alongside
third-party tools composes existing mechanisms and is not a new analyser. Two commitments shape the design.

First, \emph{normalization before comparison}. A canonical control taxonomy
mediates between native rule identifiers and benchmark ground truth, so that a
finding is credited when it identifies the control a case was authored to
violate rather than merely because the tool emitted something. Findings whose
rule identifiers the taxonomy does not cover are counted and recorded in the
released results rather than discarded, since silently treating an unmapped detection as a miss biases
the comparison toward whichever tool the taxonomy was authored around.

Second, \emph{measurement separated from assumption}. A case is admitted only by
a mechanically enforced rule: valid HCL2 referencing resource types the declared
provider defines, a non-contradictory ground-truth label, and an unambiguous
mapping onto canonical controls. The reported corpus size is the admissible count, not the
number of entries in a catalogue. A tool that is not installed in the measurement
environment is reported as such and excluded; it is never represented by an
assumed detection rate. Latency is sampled repeatedly and reported with
dispersion.

The design follows the benchmarking requirements \citet{hasselbring2021benchmarking}
sets out for software engineering research, in particular that a benchmark
publish its admission rule and its measurement configuration rather than only its
headline scores.

One outcome should be stated at the outset rather than discovered in
Section~\ref{sec:results}. Applied to the pipeline this paper evaluates, the
procedure reports a third-party scanner ahead of it on point estimates. We report
that unadjusted; it is one indication, though not proof, that the procedure does
not flatter the party that wrote it.

\subsection{Contributions}
% Three bullets, not five. The two that were cut -- the admissibility procedure
% and the execution harness -- are engineering that supports the first and third
% rather than results in their own right, and on this corpus the admissibility
% gate rejected nothing (Section~\ref{sec:admissibility}), which is not a claim
% that survives a reader who opens Table~\ref{tab:admissibility}.
\begin{itemize}
\item A finding-normalization methodology and canonical control taxonomy that
render heterogeneous IaC scanner output commensurable, with three explicit
matching criteria that separate correct attribution from incidental detection,
and a measurement of how far the choice among them moves the conclusion. On this
corpus that choice moves recall by 29 to 43 percent of the spread between the
tools compared, and moves one layer from last to second.
\item A statistical protocol for paired, small-sample scanner comparison: exact
Clopper--Pearson intervals \citep{clopper1934use}, the exact binomial form of
McNemar's test \citep{mcnemar1947note, dietterich1998approximate},
Holm--Bonferroni correction \citep{holm1979simple}, and $d_{\min}$, the smallest
discordant imbalance a comparison could have resolved. Under the correction
applied here, three of the six comparisons among the four source- and plan-level
tools could not have reached significance under any outcome, and the other three
only under a unanimous split; $d_{\min}$ is what distinguishes that from evidence
of similarity, and it is computable from the discordant count
alone.
\item An open replication package --- corpus, generator specifications,
admissibility gate, normalization engine, execution harness, and analysis code
--- in which every reported table is a generated artefact traceable to recorded
tool output, and which retains the measurements taken before the disclosed
plan-level policy correction so that the disclosure can be checked rather than taken on trust.
\end{itemize}

%------------------------------------------------------------------------------
\section{Background and Related Work}
\label{sec:related}

\subsection{Security Weaknesses in Infrastructure as Code}
\citet{rahman2019sins} established that infrastructure scripts contain recurring
security smells, first across Puppet and subsequently in Ansible and Chef
\citep{rahman2021replication}, and derived a defect taxonomy for IaC scripts
\citep{rahman2020gang}. \citet{sharma2016smell} had earlier catalogued
configuration smells over 4,621 Puppet repositories, establishing the
maintainability counterpart to the security work. Complementary studies
characterise misconfiguration in Kubernetes manifests \citep{rahman2023k8s},
Kubernetes security practices \citep{shamim2020xi}, security practice in Terraform repositories specifically
\citep{verdet2023exploring, verdet2025adoption}, and practitioner conventions
drawn from grey literature \citep{kumara2021dos}. Quality and defect-prediction
metrics for infrastructure code have also been catalogued
\citep{dallapalma2020catalog, dallapalma2022within}. This body of work motivates
automated pre-deployment validation but does not address how competing validation
tools should be compared.

\subsection{Static Analysis of IaC and Existing Comparisons}
\citet{chiari2022static} survey static analysis for IaC and note the absence of
shared evaluation infrastructure. The closest detector-side prior art is GLITCH,
which detects security smells across several IaC technologies via an intermediate
representation and is evaluated against manually annotated oracle datasets with a
tool comparison \citep{saavedra2022glitch, saavedra2023polyglot}. The contrast
with this work is a clean one: GLITCH normalizes \emph{inputs}, so that one
detector can analyse five languages, whereas IaCSecBench normalizes
\emph{outputs}, so that detectors reporting incommensurable findings can be
scored against one ground truth. Section~\ref{sec:external} states why GLITCH's
oracles are not the corpus used here and what would be required to use them. The
taxonomy such detectors are scored against is itself under revision, most
recently beyond the original seven smell categories \citep{war2025taxonomy},
which is one reason this paper versions its control set as data rather than
fixing it in prose.

Comparative evaluation of the practitioner scanners has begun to appear.
\citet{roe2026multicloud} benchmark IaC scanners across multiple cloud providers
and report per-tool detection counts; \citet{gokhale2026comparative} contrast
policy-driven evaluation with taint-aware static reasoning for IaC
misconfiguration; and \citet{kapetanidou2025k8s} evaluate the equivalent tool
class for Kubernetes. These establish that the comparison is wanted. None
addresses the two obstacles this paper targets: findings are compared in
whichever schema each tool emits, so a rule that fires for the wrong reason is
indistinguishable from one that fires for the right one, and the corpora are not
released with an admissibility procedure a reader can re-run. Our contribution is
orthogonal to these studies and could be applied to their corpora.

\citet{opdebeeck2023control} show that control- and data-flow analysis materially
changes security-smell detection in IaC \citep[see also][]{opdebeeck2022smelly}.
That result motivates the plan-level layer evaluated here: a compiled Terraform
plan is a resolved artefact in which interpolation, conditionals and expansion
have already been applied, so evaluating it sidesteps a class of analysis
difficulty instead of solving it. \citet{qiu2024semantic} pursue the
complementary route of recovering semantic checks over IaC programs directly, and
policy-violation analysis has been extended to packaged Kubernetes artefacts
\citep{minna2026helm}. Empirical work on policy-as-code adoption indicates which
controls practitioners actually enforce \citep{foalem2026pac}.

\subsection{Evaluation Methodology}
Methodologically this paper follows the static-analyser evaluation tradition
established by \citet{habib2018howmany}, who show that headline detection rates
are highly sensitive to what counts as a match, and the benchmark-construction
cautions of the OWASP Benchmark project \citep{owasp2023benchmark}. The
three-level matching criterion of Section~\ref{sec:normalization} is a direct
response to the first. Our reporting conventions are those of the static application security testing
(SAST) tool-comparison literature, which reports per-category rather than pooled
effectiveness and separates the absence of a rule from the failure of one
\citep{li2023sast, esposito2024extensive}. \citet{bohme2022reliability}
show, for fuzzer benchmarking, that a benchmark's conclusions depend on evaluation
choices a report may leave unstated, which is the argument for the version
manifest of Section~\ref{sec:environment} and for preferring interval width to
corpus size as an indicator of evidential strength. Threats to validity are
organised under the categories of \citet{wohlin2012experimentation}, and the
study is reported against the empirical standards for software engineering
research \citep{ralph2021standards, hasselbring2021benchmarking}.

\subsection{Policy as Code and Native Infrastructure Testing}
Open Policy Agent (OPA) provides general-purpose policy evaluation over structured
documents \citep{openpolicyagent2024}, and Terraform provides a native test
harness for module behaviour \citep{hashicorp2024tests}. Compliance frameworks
such as the Center for Internet Security (CIS) Amazon Web Services Foundations
Benchmark, version 3.0.0 \citep{cis2024aws}, and NIST SP~800-53
\citep{nist2020sp80053} catalogue the requirements that scanner rules are
commonly mapped to; where a benchmark control has a CIS counterpart, the
benchmark cites it. IaCSecBench
composes these mechanisms; it does not extend them.

%------------------------------------------------------------------------------
\section{Framework Architecture}
\label{sec:architecture}

IaCSecBench composes three validation layers, each addressing a distinct failure
class. Fig.~\ref{fig:pipeline} shows their arrangement and the gating relation
between them: a commit reaches a later layer only by passing every earlier one,
so the layers compose as a conjunction in execution.

Detection, however, is scored as a union, and the two are consistent. Under
gating, a case that any layer would catch is a case the pipeline catches, because
the commit is stopped at the first layer that objects; the conjunction governs
which layers get to run, not which defects get caught. Where this paper reports
the pipeline's detection effectiveness, the figure is therefore the union of
layers 1 and 3, marked as such.

Two of the three layers are scored per benchmark case. Layer 2 validates module
behaviour rather than detecting resource misconfiguration, so it has no detection
count on a corpus of misconfigured resources and is reported by suite size alone
in Section~\ref{sec:rq3}.

% Rendered vector PDF, 375x454pt (aspect 0.83). Source:
% figures/pipeline_architecture.mmd, generated from the recorded results by
% experiments/generate_figures.py.
\paperfigure{figures/pipeline_architecture.pdf}{Layered validation pipeline:
repository-edge data validation, native module testing, and policy evaluation over
compiled plan documents. A change reaches a later layer only by passing every
earlier one. Edge labels give the mean per-case latency of the scored layers; the
full figures, including dispersion, are in Table~\ref{tab:latency}. Layer 2
validates module behaviour rather than resource configuration, so it carries no
detection count on this corpus and is reported by suite size in
Section~\ref{sec:rq3}.}{fig:pipeline}

\subsection{Layer 1: Repository-Edge Data Validation}
The first layer prevents sensitive material from entering version control through
schema verification, detection of sensitive fields, and pattern-based
identification of personal data. It operates on repository content rather than
infrastructure semantics, and its characteristic failure mode is over-detection:
high-entropy strings in output declarations are frequently flagged as
credentials. Over-detection is not a cosmetic cost. Developers discount tools
whose warnings they learn to distrust, an effect documented both for static
analysis generally \citep{vassallo2019engage} and for security warnings
specifically \citep{tahaei2021notifications}, to the point that identifying which
warnings merit action is now a research area in its own right
\citep{ge2023actionable}. This layer is therefore reported as its own row in every
table of Section~\ref{sec:results} rather than pooled with the policy layers, so
its false positives stay visible.

\subsection{Layer 2: Native Module Testing}
The second layer validates module behaviour using Terraform's native test
framework \citep{hashicorp2024tests}, covering variable constraints, output
correctness, conditional resource creation, and tagging requirements without
provisioning cloud resources.

\subsection{Layer 3: Policy Evaluation over Compiled Plans}
The third layer converts a Terraform plan to its JSON representation and
evaluates policies written in Rego, OPA's policy language, against it
\citep{openpolicyagent2024}. Because the plan
is produced after expression resolution, policies observe concrete attribute
values rather than unevaluated references. Plan generation is configured so that
provider credential validation is disabled, which permits evaluation in
continuous integration without cloud account access; this is a precondition for
the layer being usable at all, and is implemented by an injected provider
override, never by relaxing a policy.

\subsection{Finding Normalization}
\label{sec:normalization}
Comparing heterogeneous scanners requires a canonical representation. A finding
is normalised to
\begin{equation}
\mathcal{F} = \langle c,\; t,\; \kappa,\; r,\; s \rangle,
\end{equation}
where $c$ identifies the benchmark case, $t$ the tool, $\kappa$ the set of
canonical controls the tool's native rule maps to, $r$ the resource address, and
$s$ the reported severity. The mapping from native rule to $\kappa$ is
many-to-many: a coarse rule may legitimately cover several controls, and such
rules are reported explicitly because they weaken attribution precision.

Detection is then evaluated at three strictness levels:
\begin{itemize}
\item \emph{control}: some finding's $\kappa$ intersects the controls the case
was authored to violate. This is the primary criterion.
\item \emph{resource}: some finding names the resource address the ground truth
identifies, irrespective of which rule fired. This isolates the
correct-resource-wrong-reason case.
\item \emph{any}: the tool emitted any finding within the case. This criterion
overstates detection and is reported to quantify how much of a tool's apparent
performance is incidental.
\end{itemize}
Reporting all three is deliberate: adopting a single permissive criterion is the
most common mechanism by which scanner benchmarks overstate detection
\citep{habib2018howmany}. Where the ground truth does not name a resource
address, the resource criterion is inapplicable, and such cases are left out of
its denominator; on this corpus that excludes the four external cases.

The three are not a single ordering. Only \emph{control} and \emph{any} nest,
because a finding that maps to the case's control is a finding. \emph{Resource}
is a different question from either, and it departs from \emph{control} in both
directions. On the 26 vulnerable cases whose ground truth names a resource
address, it credits each of Checkov, tfsec, Trivy and the plan-level layer with
one case more than \emph{control} does, and the repository-edge layer with one
fewer: on those cases \emph{control} gives 88.46, 73.08, 69.23, 61.54 and 3.85
percent and \emph{resource} 92.31, 76.92, 73.08, 65.38 and 0.00. A tool may fire on
the right resource for a different control, which \emph{resource} credits and
\emph{control} does not, or attribute the right control to a neighbouring resource
address --- the bucket rather than the encryption-configuration resource the
ground truth names, say --- which \emph{control} credits and \emph{resource} does
not. On this corpus the first accounts for five tool--case pairs and the second
for one, and the repository-edge layer's single control-level detection carries
no resource address at all. Where this paper speaks of relaxing the
criterion, and where the $\Delta$ column of Table~\ref{tab:strictness} is
computed, the movement is from \emph{control} to \emph{any} along the one axis
that does nest. \emph{Resource} is reported alongside them as a separate
diagnostic of attribution precision, not as an intermediate point on that axis.

Fig.~\ref{fig:normalization} traces a finding through the engine. The ordering
matters for the comparison's validity: normalization onto the canonical taxonomy
precedes confusion-matrix construction, so no tool is scored in its own reporting
schema and the matching criterion is applied identically to every tool.

% Rendered vector PDF, 600x486pt (aspect 1.23). Source:
% figures/normalization_workflow.mmd, generated from the recorded results.
\paperfigure{figures/normalization_workflow.pdf}{Finding-normalization engine and
evaluation pipeline: heterogeneous scanner output is mapped onto the canonical
control taxonomy before any confusion matrix is
formed.}{fig:normalization}

%------------------------------------------------------------------------------
\section{Evaluation Methodology}
\label{sec:method}

The threat classes the pipeline's controls were selected against, and the four
domains held out of scope, are given in Appendix~\ref{app:threats}. That
decomposition delimits intent; the operative statement of what this evaluation
measured is the corpus coverage of Section~\ref{sec:coverage}.

\subsection{Corpus Construction and Admissibility}
\label{sec:corpus}
\label{sec:admissibility}
The corpus separates a controlled collection, authored for this evaluation, from
an external collection of four configurations illustrating published CIS
controls. The distinction was meant to separate cases that share an author with
the policy set, and so cannot confirm detection capability independently, from
cases that do not. The replication package does not record an upstream source for
the external configurations, however, and they entered the repository in a commit
by the investigator, so we do not rely on them for independence
(Section~\ref{sec:external}).

Admissibility is enforced mechanically. A case is admitted only if it satisfies
all of the following:
\begin{enumerate}
\item it declares at least one resource or module block;
\item its ground-truth label is present and internally consistent, where a case
declaring contradictory labels is rejected rather than resolved by precedence;
\item it resolves to exactly one interpretation in the canonical taxonomy, since
a CIS section number may span several distinct controls and inference cannot
choose between them;
\item \texttt{terraform init} and \texttt{terraform validate} succeed, which is
the only check that detects a case referencing a resource type the provider does
not define.
\end{enumerate}

\resulttable{corpus}

Table~\ref{tab:admissibility} reports the outcome, per collection. The admissible
count is the corpus size used throughout; catalogue entries without corresponding
configuration are not cases and are excluded. The two collections diverge by very
different factors, which is why the table is split: the controlled catalogue
declares 345 entries against 52 configurations, and the external catalogue 175
against 4.

On the released corpus the gate rejected nothing: every configuration present on
disk was admitted. We report that rather than leaving the reader to infer a
filtering effect from the presence of a filter. What the procedure buys here is
not a rejection count but a stated and re-runnable admission rule, which is the
property \citet{hasselbring2021benchmarking} asks a benchmark to publish; its
value as a filter is untested on a corpus this clean and would be established
only by a corpus that fails it.

Both gaps require explanation, and they have different causes. The external
manifests declare 50 CIS, 75 SecureFlag and 50 Terraform Registry entries, of
which only four CIS configurations exist and the other two collections have none; a manifest entry
without a configuration is a citation rather than a case, and counting such
entries would inflate the external collection by more than a factor of forty
while contributing no measurement. The controlled catalogue is the residue of a
larger planned corpus: none of its 345 entries has a configuration on disk, and
the 52 controlled cases were instead generated from per-control specifications
under identifiers of their own. We
report the declared count rather than pruning the catalogue to match, because the
gap is itself the datum that motivates reporting admissible rather than
catalogued size, and because the catalogue indicates where the corpus can be
extended. No reported figure derives from a declared-but-absent entry.

\subsection{Control Coverage}
\label{sec:coverage}
The canonical taxonomy defines 26 controls, and every one is now exercised by at
least one admissible case. An earlier revision of this corpus exercised 22: four
controls, one for API-gateway authorization and three for Kubernetes workloads,
were declared in the taxonomy with nothing run against them. A control in that
state is invisible in the results rather than reported as a gap, because an absent
case produces no row, so the four were specified and generated.

Closing them established a distinction the confusion matrix cannot express.
Neither tfsec nor Trivy emits any finding
on any \texttt{kubernetes\_*} Terraform resource --- not a different finding, none
at all, verified on both \texttt{kubernetes\_pod} and
\texttt{kubernetes\_pod\_security\_policy} and on the compliant variants as well as
the vulnerable ones. Only Checkov inspects them. Their non-detection on those three
controls therefore measures the scope of what each tool reads, not its accuracy,
while \S\ref{sec:rq2} scores it identically to a missed misconfiguration. The
plan-level layer's misses on the same cases have a different cause: its policy
set implements no rule for the three controls, so they are a coverage gap in the
reference implementation and are scored as misses. The control map records tfsec
and Trivy, and only those two, as inapplicable for those controls and the
analysis emits a caveat naming them, but the tables cannot separate the two
failures, so we report the effect directly: including the Kubernetes controls
moves control-level recall by $+1.1$ points for Checkov, $-8.5$ for tfsec, $-8.1$
for Trivy and $-7.4$ for the plan-level layer. Treating tfsec's and Trivy's
Kubernetes cases as inapplicable rather than missed would raise their recall to
85.19 and 81.48 but leave the best-to-worst spread at 23.33 points, because its
endpoints are Checkov and the plan-level layer (\S\ref{sec:rq1}).

Auditing the taxonomy against the installed tool catalogues found three defect
classes in the mapping itself, all corrected. Seven controls named plan-level
policy rules the policy set does not implement. Two named Checkov identifiers
absent from the installed catalogue entirely, and one mapped a
resource-\emph{requests} check to a resource-\emph{limits} control, which are
distinct Kubernetes concepts. None of these can produce a spurious detection,
because nothing emits the identifier; each instead inflates the apparent coverage
of the layer it is attributed to, and in the first case that layer is the
reference implementation. That is the direction of error that flatters the
contribution under evaluation, which is why the audit was mechanical, and why the
resulting invariant, that every plan-level identifier must exist in the policy
source, is now enforced by a test. A later audit found one defect that could have
produced a spurious detection: tfsec's and Trivy's Lambda tracing rules were
mapped to the function-encryption control. Neither fires on that control's cases,
so no classification changed, and both were removed; tfsec and Trivy now have no
rule for that control, so their miss on it is an absent rule rather than a failed
one.

The correction left two controls with no implementing rule in any tool, and so
detectable by none of them. A control in that state scores every case citing it
as a miss for every tool simultaneously, which reads as unanimous tool failure
rather than as absent coverage; both were removed, and neither was exercised, so
no reported figure changes. Two controls retain an unverified flag, down from
eight. The eight were retained on the grounds that this environment could not
corroborate a tfsec long identifier: the recorded tfsec release provides no
check-listing command,
and its \texttt{--filter-results} option returns nothing for an unknown identifier
and for a real but untriggered one alike. Purpose-built cases settle the question
that the tool's interface cannot. An identifier that fires on the vulnerable member
of a pair and stays silent on the compliant one is corroborated by observation
regardless of what the catalogue can be made to list, and on that basis three of
the eight were confirmed against emitted output and their evidence recorded
alongside the mapping. Three more, the Kubernetes controls, were resolved rather
than confirmed: tfsec emits nothing on Kubernetes resources, so its identifiers
there were moved out of scope and tfsec and Trivy recorded as inapplicable, with
Checkov's firing recorded as the evidence for each control.

The two survivors are comparator gaps rather than mapping errors, and we state what
was observed instead. For a compute instance assigned a public address, Checkov
fires but tfsec emits nothing for \texttt{associate\_public\_ip\_address} on an
instance; it flags only the subnet-level form when
\texttt{map\_public\_ip\_on\_launch} is set, which this case does not set. For
credential material embedded in configuration, no comparator emits a finding
mapped to that control, although all three report other findings on the same
resource: Checkov's
secret checks belong to its secrets framework and the harness invokes it with
\texttt{--framework terraform}, so they cannot fire, and only the reference
pipeline's repository-edge layer detects the case. That second asymmetry favours
the contribution under evaluation, which is why it is stated here rather than left
in a mapping file. Identifiers are retained rather than deleted, because a wrong
identifier is inert and the only possible effect of keeping it is to credit a
comparator if it turns out to be real. Where the two options differ we take the one
that cannot flatter the reference.

\subsection{Ground-Truth Labelling}
\label{sec:labelling}
Labels are derived mechanically rather than by human rating. Each controlled case
is emitted by a generator from a per-control specification that declares the
intended violation, and the label follows from which member of the
vulnerable/compliant pair is being written. The generator, the specification, and
the emitted cases are all in the replication package, so any reader can check
that a case carries the label its specification declares.

Because labels are a deterministic function of the specifications, rater
disagreement is impossible by construction, and the threat that replaces it is
\emph{specification error}: a specification that misreads its originating control
produces a wrong label reproduced identically across the whole pair. An agreement
coefficient computed over the generator's own output would measure its
determinism, not concurrence, so we do not report one.

We instead ran a blind second labelling pass against the threat that does apply.
It is an \emph{automated consistency audit}, not a second rating in the
inter-rater sense: \textbf{the second labeller is a large language model, not a
second human investigator}. We name it that way throughout because the word
\emph{rater} imports a reliability claim this design cannot support, and the
distinction governs how the agreement below may be read.

The identity and version of the auditing model, its access route and decoding
parameters, and the exact instruction it was given were not recorded when the
pass was run, and we do not reconstruct them after the fact. What is recorded is
the blinding procedure described below and the label and reason returned for
each case, in \path{benchmark/labelling/independent_relabelling.json}. No case's
labelling carried another case's label or reason by construction, but the cases
were labelled one at a time within a single session, so earlier cases' labels and
reasons may have been in context when later ones were labelled, and the record
cannot exclude it. Five cases are a more direct exception, discussed at the end
of this section: the model had already seen their configurations and recorded
labels. The omission matters: the agreement reported below is a
property of one model under one instruction, and with neither identified the
labelling cannot be repeated. The per-case labels are recorded, so the agreement
figure can be recomputed from them, but it cannot be re-obtained from the
instrument. We therefore report
it as a record of what was observed, not as a reproducible measurement, and the
single disagreement it surfaced stands on its own analysis below rather than on
the audit's authority. The model was used as a measuring
instrument in this role; it is not an author, and responsibility for how its
labels are used here rests with the author. The other uses of AI assistance in
this study are described in Section~\ref{sec:environment}.

For each admissible case a review view was constructed containing only the
Terraform configuration and the canonical control under test, with every comment
line removed, since the generator writes the expected label and its rationale
into each file header, and with the annotation files and the case identifier
withheld, the latter because its suffix names the class. Cases were presented in
an order keyed on a hash of the case identifier, and the views were screened for
residual label vocabulary. The recorded view contained only the configuration and
the control text; the instruction that accompanied it was not recorded. A label
and a reason were recorded for each case before any comparison.

Agreement is reported before reconciliation, since resolving disagreements
guarantees consensus: 47 of 48 cases agreed, giving raw agreement of 97.92\% and
Cohen's $\kappa = 0.958$ against a chance-agreement baseline of 50.17\%. We report
$\kappa$ as a consistency figure for one model under one prompt and not as an
inter-rater reliability coefficient, for which this design supplies no evidence. With
the class balance of the audited set, 26 vulnerable against 22 compliant, expected
agreement is close to one half, so $\kappa$ here carries little information beyond
the raw count; the $2 \times 2$ agreement table is in the replication package.

The audited set is 48 of the 56 admissible cases, and the eight it excludes are
the four pairs added to close the coverage gap of \S\ref{sec:coverage}. The
relabelling pass was run once, blind, against the corpus as it then stood; the
later cases have not been through it, and running the model again would not extend
the same audit, because the original model and instruction are unrecorded and a
new run would be a different instrument. We therefore report the
audit on the set it covered rather than extrapolating its agreement rate to the
whole corpus, and note that the eight unaudited cases are precisely the ones whose
labels were written last and reviewed least.

The single disagreement is the kind of defect this pass exists to surface. A
bucket configured with SSE-S3, server-side encryption with Amazon S3 managed keys
(\texttt{AES256}), was labelled violating under a
control titled \emph{lacks server-side encryption} and citing CIS Amazon Web
Services 2.1.1. That is a requirement of version 1.4.0 of the benchmark,
\emph{Ensure all S3 buckets employ encryption-at-rest} \citep{cis2021aws}, which
version 3.0.0 removed. The configuration does not lack server-side encryption,
and that requirement is satisfied by SSE-S3; read against its stated text the
case is compliant.
The generator's rationale reveals the intent, namely encryption \emph{other than}
a customer-managed key, so the label encoded a requirement the control text never
stated. That stricter requirement is also what the reference policy set enforces,
which makes this a concrete instance of the construct-validity threat of
Section~\ref{sec:construct} rather than a hypothetical one: the ground truth
agreed with the reference implementation rather than with the control it cited.

We corrected the control text rather than the label, because requiring a
customer-managed key is a defensible control and is what the corpus, the policy
set, and the customer-key rules of all three source-level scanners test;
relabelling instead would have left the pair testing nothing. The control is
retitled to state the key requirement and records that it is stricter than the
version 1.4.0 requirement it derives from; it cites no version 3.0.0 requirement,
because there is none. No confusion-matrix entry changed. A regression test now
pins the citation in the taxonomy to the one in the generator specification so
the two cannot drift apart silently.

A later audit of every citation against version 3.0.0 \citep{cis2024aws} found the
taxonomy mixing its numbering with that of version 1.4.0, and several citations
naming unrelated requirements. All are now given against version 3.0.0; the five
controls whose citations had no counterpart there cite none, with the reason
recorded in the control map. The four external configurations carry labels of
their own that follow no version's numbering, and their metadata records the
corresponding version 3.0.0 requirement. No label or classification depends on a
citation number, so no reported result changed.

Three properties bound what this pass establishes, and we state each at the
strength the record supports. The audit reads the control texts from the same
public documentation the specifications were written from, so the two readings
are not fully independent and agreement is biased upward accordingly. The model
was run in the same assistant session as earlier debugging, during which it had
seen the configurations and recorded labels of five cases, so their agreement is
not blind; excluding all five leaves 42 of 43
agreeing, with the disagreement outside that set. And the blinding was not
perfect: one external case declares \texttt{resource "aws\_cloudtrail"
"insecure\_trail"}, a resource name that hints at the label. The screen for
residual label vocabulary flagged it, and it was left unmodified rather than renamed, so that the released
case is the one that was audited; the case is identified in the replication
package so its contribution to the agreement figure can be discounted. The full
per-case labels, reasons, and these caveats are in the replication package.

Two further checks bear on specification error. The admissibility procedure of
Section~\ref{sec:admissibility} rejects any case whose label is absent,
internally contradictory, or inferred to resolve to more than one control in the
taxonomy; a case may still declare several controls explicitly, as one external
CIS case does.
And the audit of the control map reported in Section~\ref{sec:coverage} found
three classes of mapping defect by mechanical comparison against the installed
tool catalogues. Neither check can establish that a specification reads its
control correctly; that remains an unmitigated threat, and independent
relabelling by a second investigator is the obvious next step.

\subsection{Research Questions}
The matching criterion is the primary object of study, and the research questions
are ordered accordingly.
\begin{itemize}
\item \textbf{RQ1}: How much does the matching criterion change apparent
detection, relative to the differences between the tools being compared?
\item \textbf{RQ2}: What detection effectiveness do the evaluated tools achieve
on the admissible corpus, how precisely is that effectiveness estimated, and what
difference could the comparison have detected?
\item \textbf{RQ3}: What execution latency does each layer contribute in
continuous integration, and what is the size of the native validation suite
required to express the corresponding assertions?
\item \textbf{RQ4}: Which layers account for which detections?
\item \textbf{RQ5}: Does effectiveness persist on independently sourced
configurations?
\end{itemize}

The inferential statistics that answer RQ2 are there to characterise uncertainty
and to bound what a comparison at this corpus size could resolve. They are not
there to establish that one tool is superior, and no claim of superiority or of
equivalence is drawn from them anywhere in this paper.

\subsection{Measurement Environment and Tool Versions}
\label{sec:environment}
Tool versions are resolved at execution time and recorded in
\path{results/run_manifest.json}, together with the platform, processor, and
repository commit; for the reported run that commit does not identify the
measured tree, as Section~\ref{sec:construct} sets out. Versions are not restated in prose, because a hand-copied
version string is the first thing to drift; the manifest is authoritative
\citep{bohme2022reliability}.

Tools not installed in the measurement environment are recorded as such and
omitted from every table. No tool is assigned an assumed detection rate or an
assumed latency.

AI coding assistants were used in building the study: to write parts of the
harness and analysis code, the plan-level policy set and the benchmark cases, and
to help debug the harness and inspect results, and an AI assistant was used to
draft revisions of manuscript passages under the author's direction. Which
assistant products were used is stated in the Statements and Declarations. None of these
uses is an instrument of the detection measurements: every reported detection and
latency figure is the recorded output of the tools under test, scored by the
released harness, and every artefact the assistants contributed to is in the
replication package for inspection. The one instrument role a language model did
play, as second labeller in the consistency audit, is described in
Section~\ref{sec:labelling}. What their use in writing both the
policies and the cases implies for construct validity is discussed in
Section~\ref{sec:construct}.

\subsection{Tool Selection}
\label{sec:toolsel}
The comparison covers Checkov, tfsec, Trivy, and standalone policy evaluation
over plan documents. Selection was governed by two criteria: the tool must parse
HCL2 or Terraform plan JSON without requiring cloud provider API access, and the
set must span distinct analysis paradigms, namely syntax-tree analysis, lexical
scanning, and policy evaluation. KICS, Regula, and provider-native configuration
services are excluded under the first criterion or as duplicating a represented
paradigm.

Trivy requires separate comment, because it is in the set for a reason unrelated
to paradigm coverage. Aqua Security folded tfsec's engine and rule set into
Trivy, and Trivy is the maintained product; tfsec is no longer developed. A
comparison reporting tfsec alone would be characterising a scanner practitioners
are advised off, so both are measured. But the two are not independent
implementations. Their rule identifiers are the same Aqua Vulnerability Database
numbers under different spellings (tfsec reports \texttt{AVD-AWS-0086} where
Trivy reports \texttt{AWS-0086}), and we exploit that correspondence to build
Trivy's entries in the control taxonomy: each is derived mechanically from the
identifier pair tfsec itself emits, rather than chosen by observing which rules
Trivy happens to fire on a case, which would fit the taxonomy to the tool. A
regression test re-derives the crosswalk and fails if the two disagree.

The consequence bounds what the enlarged tool set establishes. Three third-party
scanners are measured, but only \emph{two} independent third-party rule sets:
Checkov's, and the tfsec/Trivy lineage. Agreement between tfsec and Trivy is
therefore close to a consistency check on that shared lineage across two engine
generations rather than corroboration by a second party, and we do not report it
as the latter. Where the two diverge, the difference is attributable to engine and
rule evolution between the products, and is informative for that reason.

Terrascan is absent for a reason that is not a selection criterion, and we state
it plainly: no successful Terrascan execution was ever recorded against this
corpus, so there is nothing to report. Its paradigm, Rego policy evaluation, is
represented in the tool set by the plan-level layer, which bounds what the
omission costs without eliminating it. Policy evaluation is therefore represented
only over compiled plan documents, and no source-level policy engine appears in
the tool set. A source-level Rego engine could differ from a plan-level one on
configurations whose violations are resolved during planning, and this evaluation
does not measure that difference. We make no claim of any kind about Terrascan's
detection behaviour.

%------------------------------------------------------------------------------
\section{Results}
\label{sec:results}

All tables in this section are generated by \path{evaluation/analyze.py} from
recorded raw tool output, regenerated by \path{experiments/run_baselines.sh},
and not edited by hand. In them, \emph{IaCSecBench L1} is the repository-edge
layer and \emph{OPA (plan-level)} is the plan-level layer, which evaluates the
author's policy set.

\subsection{RQ1: Sensitivity to the Matching Criterion}
\label{sec:rq1}
\resulttable{strictness}

Table~\ref{tab:strictness} reports recall under each criterion, with the
control-to-permissive spread $\Delta$ given per tool, in cases as well as in
points, and with an exact interval on the case count. Two readings of that column
matter, and they differ in magnitude by an order of magnitude.

Among the three source-level scanners and the plan-level layer, $\Delta$ ranges
from 6.67 to 10.00 percentage points. The spread between those same four tools at
the control criterion is 23.33 points, from Checkov's 90.00 to the plan-level
layer's 66.67, which is more than twice the criterion effect.

That comparison needs one qualification, which \S\ref{sec:coverage} sets out.
Three of the 26 controls govern Kubernetes workloads. tfsec and Trivy read none of
the resources those controls are expressed in, so their zero on those cases is a
statement about scope rather than detection; the plan-level layer's zero on the
same cases is a gap in its own policy set and a genuine miss. Treating tfsec's
and Trivy's Kubernetes cases as inapplicable leaves the spread at 23.33 points,
because its endpoints are Checkov and the plan-level layer. Excluding the three
controls altogether would narrow it to 14.81, but only by removing the reference
implementation's own missing rules, which flatters the reference, so we do not
treat that figure as scope-controlled. On the point estimates the criterion effect
is therefore 29 to 43 percent of the between-tool spread, and the spread is 2.3
to 3.5 times the criterion effect; with intervals as wide as those below, this corpus
does not establish that ordering beyond the point estimates.

Both effects remain small in cases: two or three of the 30 vulnerable cases move
for each of the four source- and plan-level tools when the criterion is relaxed. The intervals in
Table~\ref{tab:strictness} are correspondingly wide, so what this corpus
establishes is that the criterion effect is not negligible beside the
between-tool effect, not the size of either. We state the claim at that strength
and no further.

Rank order is where the criterion's effect is uneven, and it is worth being
precise about where it does and does not disturb it. Among the three source-level scanners and the
plan-level layer the criterion reorders nothing: Checkov leads at both criteria,
the plan-level layer trails at both, and relaxation only collapses tfsec and
Trivy into a tie at 83.33. The
reordering is entirely the repository-edge layer's, and it is large. That layer's
$\Delta$ is 90.00 points, from 3.33 at the control criterion to 93.33 under the
permissive one, which moves it from last of the five to second. It detects
hardcoded credentials and comparable textual patterns, so on a corpus dominated
by attribute-level controls it identifies the seeded violation in exactly one
case while emitting some finding in almost all of them. The ratio between its two
figures is close to twenty-eight, and we report the difference in points rather
than that ratio: with one true positive in the denominator the ratio is an
artefact of a small count, whereas the 90-point difference is the quantity that
would mislead a reader of a single-criterion report.

The two readings support different claims, and only the second is strong here. A
single-criterion report does not necessarily misstate the \emph{ranking} of tools
working at the same abstraction level; on this corpus it does not. It can
misstate the standing of a tool whose findings are of a different kind by most of
the available range, and it gives a reader no way to tell which of the two
situations they are in. That is the argument for reporting the criterion, and it
does not require that the criterion have reordered anything.

In answer to RQ1, relaxing the criterion from \emph{control} to \emph{any} moves
recall by 6.67 to 10.00 points for the source- and plan-level tools, 29 to 43
percent of the 23.33-point spread between them, and by 90.00 points for the
repository-edge layer, which it moves from last to second.

\subsection{RQ2: Detection Effectiveness and Resolvable Differences}
\label{sec:rq2}
\resulttable{performance}
\resulttable{rates}

Table~\ref{tab:performance} reports confusion-matrix counts and
Table~\ref{tab:rates} the corresponding rates with exact Clopper--Pearson
intervals. Intervals accompany every proportion because a point estimate on a
corpus of this size is compatible with a wide range of true values; Checkov and
the plan-level layer, attaining 100.00\% precision, carry lower confidence bounds
of 87.23 and 83.16 respectively, and reporting the point estimate alone would misstate what has
been established.

The Clopper--Pearson form is chosen for its coverage guarantee rather than for
interval width. Its actual coverage is never below the nominal level, whereas the
score and adjusted-Wald intervals attain nominal coverage only on average and
undercover for proportions near the boundary \citep{agresti1998approximate},
which is where most of these estimates sit; Agresti and Coull draw the opposite
recommendation from the same observation, preferring the approximate intervals
for their width. The cost is conservatism, so every
interval in Table~\ref{tab:rates} is wider than a score interval on the same
counts; for a claim about detection capability that is the safe direction of
error, and the harness computes the score interval as well, so a reader
preferring it can obtain it from the released results.

One property of the precision intervals should be read with care. Recall's
denominator is the 30 vulnerable cases, fixed by the corpus design, so a binomial
interval applies directly. Precision's denominator is the number of findings the
tool emitted, which is a realised outcome rather than a design parameter, so its
interval is conditional on the observed prediction count and is not a marginal
statement about the tool.

Specificity is estimable here because the corpus carries 26 ground-truth
compliant cases, and it is reported for that reason: a detector's cost to its
users is carried as much by its false positives as by its misses, and a table
reporting only recall and precision would have described the half of the
behaviour that flatters a detector. The Matthews correlation coefficient (MCC)
\citep{matthews1975comparison} is likewise estimable and is given in
Table~\ref{tab:performance}. What the compliant baseline does not yet support is
a \emph{precise} false-positive estimate: with 26 negatives the specificity
intervals span more than thirteen points even where no false positive occurred.

\resulttable{mcnemar}
\resulttable{allpairs}

Tables~\ref{tab:mcnemar} and \ref{tab:allpairs} report pairwise comparisons.
Table~\ref{tab:allpairs} is the confirmatory analysis: it covers all ten
unordered pairs with Holm--Bonferroni applied across the whole set, so the
comparison a practitioner most wants, between two third-party scanners, is made,
and the pipeline evaluated here does not sit on one side of every test.
Table~\ref{tab:mcnemar} retains the against-one-reference form for continuity
with the way this literature usually reports, and its weaker correction reflects
its smaller family; where the two disagree, the all-pairs adjustment governs.

The discordant count $b$ includes every case in which the first tool is correct
and the second is not, counting spurious findings as well as missed violations;
restricting $b$ to missed violations understates disagreement. The exact binomial
$p$-value is reported in preference to the chi-square approximation, which is
unreliable when a discordant cell is small; on these counts, where the smallest
discordant cell is one, the approximation is not merely conservative but invalid.
Test selection follows established guidance for software engineering experiments
\citep{arcuri2014hitchhiker}. Odds ratios carry a Haldane--Anscombe correction,
so a zero discordant cell yields a finite estimate; an unbounded ratio is a
degenerate cell rather than a large effect, and Cohen's $g$ is reported as the
effect size.

No pairwise difference among the three third-party scanners and the plan-level
layer is significant after correction. The $d_{\min}$ column states what each
comparison could have detected, and it is the more informative result. For
Checkov against the plan-level layer, eleven discordant cases mean that only a
split of ten to one or more extreme would have reached $\alpha = 0.05$; the
observed split was nine to two. Checkov against tfsec and against Trivy each have
eight discordant cases, and the plan-level layer against tfsec and against Trivy
each have seven, so each of those four comparisons could have reached
significance only if every discordant case had fallen the same way. For tfsec
against Trivy the discordant count is two, at which \emph{no} split reaches
significance: the smallest attainable two-sided $p$-value at two discordant pairs
is $0.5$. That comparison had no power at all, and the other five had power only
against near-unanimous disagreement, so their non-significance carries little
information about the tools.

Those thresholds are uncorrected, and the correction that governs is stricter.
The four repository-edge comparisons occupy the first four Holm ranks, so any
other pair needs $p \le 0.05/6 \approx 0.0083$ to survive. Holding the other
comparisons as observed, the plan-level layer against tfsec and against Trivy
could not have survived under any split of their seven discordant cases, whose
smallest attainable $p$ is $0.0156$, nor could tfsec against Trivy; the remaining
three pairs could have survived only under a unanimous split.
Table~\ref{tab:allpairs} reports $d_{\min}$ at the uncorrected level.

The one significant family is the repository-edge layer against each of the
other four, which reflects that it addresses a different class of defect
entirely; that is a statement about scope, not quality. The interval widths in
Table~\ref{tab:rates} and the $d_{\min}$ column of Table~\ref{tab:allpairs}
together are what this corpus establishes about relative effectiveness among the
scanners, which is little, in a form a reader can act on.

In answer to RQ2, control-level recall among the four source- and plan-level tools
ranges from 66.67 for the plan-level layer to 90.00 for Checkov, with 95 percent
intervals between 24 and 36 points wide, and no difference among them survives
correction.

\subsubsection{Divergence within the shared rule lineage}
The tfsec/Trivy pair (Section~\ref{sec:toolsel}) permits one measurement the rest
of the comparison cannot make: how much a rule set changes as it is carried from
a retired product into its successor. The two agree on 54 of 56 cases. Both
disagreements concern the same control, identity and access management (IAM)
policy wildcards: tfsec reports
\texttt{aws-iam-no-policy-wildcards} on both members of the pair, and Trivy emits
its counterpart \texttt{AWS-0057} on neither, not on these two cases and nowhere
in the corpus.

The divergence is not a straightforward regression. tfsec's rule fires on the
violating case, earning a true positive, \emph{and} on the compliant case,
earning a false positive; Trivy's silence forfeits both. tfsec therefore records
the higher recall and Trivy the higher precision, from a single rule behaving
differently. We verified that Trivy's miss is genuine and not an artefact of our
taxonomy: the one unmapped finding Trivy emitted on that case is an S3
versioning rule, unrelated to IAM policy wildcards, so no mapping gap is
concealing a detection.

The remainder of the two tools' output corroborates the shared provenance
quantitatively. Every rule Trivy emitted has a tfsec counterpart reported on an
identical number of cases. The converse fails for four rules, which tfsec emits
and Trivy emits nowhere in the corpus: the IAM policy-wildcard rule above; a
CloudFront TLS-policy rule the taxonomy does not cover; and two mapped rules, S3
bucket encryption (\texttt{AWS-0088}, on 14 cases) and public database access
(\texttt{AWS-0082}, on one), neither of which changes a classification, since
the two tools disagree only on the IAM pair. Two products separated by a change of
engine thus differ, on this corpus, by four rules, in every case by the successor
being silent where the predecessor fires. Section~\ref{sec:discussion}
takes up what that implies for practitioners migrating between them.

\subsection{RQ3: Execution Cost and Suite Size}
\label{sec:rq3}
\resulttable{latency}

Table~\ref{tab:latency} reports per-case latency as a mean with standard
deviation over 168 samples per tool, three repetitions across the 56 cases.
Single-run figures are not reported, because a single sample carries no estimate
of its own dispersion.

On engineering effort we report a measurement of the native suite alone and
withdraw the comparative claim. The validation suite comprises 11
\texttt{.tftest.hcl} files totalling 190 physical lines, of which 153 are
non-blank and 137 are both non-blank and non-comment, expressing 11 \texttt{run}
blocks and 26 \texttt{assert} blocks. We give both line counts because they
differ by the 16 comment lines and one of them is the figure a reader would
compute; quoting the larger under the stricter description would inflate the
suite by 12\%. Counts were obtained by line-oriented pattern matching over the
suite sources, which are in the replication package so the count can be
reproduced or disputed.

We report no code-volume ratio against an external test framework. That would
require an independently authored equivalent suite, and constructing one
ourselves would place both sides of the comparison in the hands of the party with
an interest in the outcome. The claim that native testing reduces test-suite
volume is therefore outside this paper's scope. The native suite does require no
language runtime, dependency manifest, or build step beyond the Terraform binary
already needed to plan the configuration --- a difference in toolchain surface,
not a measured difference in effort.

In answer to RQ3, per-case latency ranges from 0.15\,ms for the repository-edge
layer to about 1.9\,s for Checkov (Table~\ref{tab:latency}), the plan-level
figure being a lower bound that excludes plan generation, and the native suite
expresses its assertions in 26 \texttt{assert} blocks across 11 files.

\subsection{RQ4: Layer Attribution}
\resulttable{layers}

Table~\ref{tab:layers} attributes detections to the layer that produced them.
Each layer is scored by the same normalization and matching procedure applied to
the external scanners, so a layer is credited only when its finding maps to the
control the case seeds; no layer receives credit by construction.

The table reports pairwise intersections and the union alongside the per-layer
counts, because the per-layer counts alone cannot distinguish a complementary
pipeline from a redundant one. Reported additively, a case detected by every
layer would be counted repeatedly and the composed pipeline would appear to
achieve more than it does. The union is therefore the only figure that
characterises the pipeline as a whole, and any excess of the summed per-layer
counts over the union is redundancy rather than additional coverage.

Two structural constraints bound what this attribution can show, and both
suppress apparent complementarity. The repository-edge layer is insensitive to
every control expressed as a resource attribute, so on this corpus its
contribution is necessarily small, as Section~\ref{sec:rq1} already showed from
the other direction. The plan-level layer can only be credited on cases for which
a plan was produced; cases whose planning failed would be excluded from its
denominator rather than recorded as misses. On this corpus no case was excluded
on that ground, so the layer's denominator is the full 30 vulnerable cases and
the constraint is stated as a property of the attribution rule rather than as an
adjustment applied here. The native module-testing layer
validates the project's own infrastructure and is not exercised per benchmark
case, so it is excluded from this table and reported under RQ3; its omission is a
scope boundary, not a null result.

In answer to RQ4, the plan-level layer accounts for 20 of the scored pipeline's 21
detections and the repository-edge layer for the remaining one, with no case
detected by both.

\subsection{RQ5: External Generalization}
\label{sec:external}

Two subsets lie outside the controlled corpus, and neither supports an
independent effectiveness estimate. The first is four labelled configurations
illustrating CIS Amazon Web Services controls (Table~\ref{tab:admissibility}),
whose labels derive from the control each illustrates; they enter the admissible
corpus on the same terms as the internal cases and are validated by the same
procedure. The replication package does not record an upstream source for them,
and they entered the repository in a commit by the investigator, so we do not
treat them as independent of the investigator.

One defect in how they were stored is worth reporting, because its signature is
invisible in aggregate. A case is scanned as a directory rather than as a file,
and these four configurations initially shared one directory. Every tool
therefore received identical input for all four cases and returned a
byte-identical finding set for each, while the scoring stage compared that one
result against four different controls. Four rows entered each confusion matrix
where one observation existed. This does not add noise so much as misstate
precision: interval width is driven by the number of independent observations, so
duplicated cases narrow every interval that includes them without contributing
evidence, and nothing in the aggregate output distinguishes this from four
genuinely independent cases that happen to agree. Each configuration now occupies
its own directory, and the corpus loader refuses to load a corpus in which two
cases share one. The same trap applies to any evaluation whose unit of analysis is
a directory: case isolation has to be enforced rather than assumed.

Four cases are too few for a per-subset recall estimate, whose interval would
span most of the unit interval, so we report them within the combined admissible
corpus.

\resulttable{external}
\resulttable{external_agreement}

The second subset is 25 Terraform repositories published by Amazon Web Services
under the MIT-0 licence, 15,226 lines in all, each pinned to an exact commit, with
the selection criteria recorded in the replication package. These are the only
configurations here whose third-party provenance is recorded, but they carry no
ground-truth label, so they yield no detection figure and nothing measured on
them enters a confusion matrix: Table~\ref{tab:external} counts alerts instead.
They bear on two points made elsewhere. The share of alerts whose native rule has
no entry in the control map is 73 percent for Checkov, 19 for tfsec and 35 for
Trivy, against 72, 25 and 27 percent on the labelled corpus. Checkov's share is nearly
the same on both, tfsec's is lower on production-shaped code and Trivy's higher,
so the labelled corpus does not overstate how much of a tool's output a
control-level comparison can score. And the tools agree little
(Table~\ref{tab:external-agreement}): by flagged resource the Jaccard index is
0.21 for Checkov against tfsec, 0.42 against Trivy, and 0.49 between tfsec and
Trivy, whose shared lineage (Section~\ref{sec:toolsel}) makes the last not an
independent comparison. Agreement is not correctness, since two tools may agree
and both be wrong.

In answer to RQ5, this study cannot say whether effectiveness persists on
independently sourced configurations: the labelled external cases are too few and
of unrecorded provenance, and the independently sourced repositories are
unlabelled.

We make no claim that effectiveness measured here transfers to production
infrastructure. Establishing that requires a labelled corpus of real-world
configurations, and the position of this literature deserves to be stated
precisely rather than as an absence. Manually annotated oracle datasets of
real-world IaC scripts do exist: GLITCH releases three, for Ansible, Chef and
Puppet \citep{saavedra2022glitch}, and later versions of the tool add Terraform
support \citep{saavedra2023polyglot}. They are not the corpus used here, first
because none of the oracles is a Terraform corpus, and second because their labels are security-smell categories in the sense of
\citet{rahman2019sins}, whereas this evaluation's unit of ground truth is a CIS
control identifier, and the two taxonomies are not in correspondence: a single
smell category spans several CIS controls and several controls have no smell
counterpart. Scoring the tools evaluated here against those oracles would require
a mapping between the two taxonomies that does not exist and whose construction
is a research contribution in itself. We identify that mapping, rather than the
absence of labelled real-world data, as the principal obstacle to
field-representative evaluation in this area, and return to it in
Section~\ref{sec:conclusion}.

%------------------------------------------------------------------------------
\section{Discussion}
\label{sec:discussion}

Five findings from building the harness bear on how IaC security tooling should
be evaluated.

First, the matching criterion is a consequential choice in its own right. Because
native rule identifiers carry no shared semantics, any comparison implicitly
chooses what counts as a correct detection, and on this corpus that choice moves
recall by 29 to 43 percent of the spread between the tools compared, and by most of the available range for one layer
(Section~\ref{sec:rq1}). \citet{habib2018howmany} established the sensitivity for
analysers of one language, which share a vocabulary of defect kinds even where
they disagree about instances. The IaC case is the harder one: the tools compared
here work at different abstraction levels, from lexical scanning to policy
evaluation over resolved plans, and there is no shared rule vocabulary in which
to state what a correct detection would be. The criterion is therefore not a
tuning parameter of the comparison but a construction of its dependent variable,
and it belongs in a benchmark's reported method alongside its corpus size. The
per-category reporting discipline the SAST literature arrived at is the same
response to the same problem \citep{li2023sast}.

Second, taxonomy coverage is a bias mechanism. When a canonical taxonomy is
authored around one tool's rule set, findings that other tools emit under
identifiers the taxonomy lacks are scored as misses. The effect, where it occurs,
is systematic and favours the authoring tool. Auditing observed rule identifiers against the
taxonomy after execution, and reporting the unmapped remainder, is the minimum
control. On this corpus 367 findings fell outside the control map, 255 of
Checkov's, 57 of tfsec's and 55 of Trivy's, but only 11 of them sit on vulnerable
cases the emitting tool missed (8, 1 and 2 respectively). They fall on five
tool--case pairs, so at most five misses could reflect a gap in the map rather
than in the tool. During this work the audit identified rules from comparator tools that
the initial taxonomy omitted and that would otherwise have been scored as
failures.

A third observation concerns statistical reporting rather than IaC. On corpora of
this size a non-significant paired test is routinely read as evidence of
similarity, yet under the correction applied here three of the six comparisons
among the four source- and plan-level tools could not have been significant
under any outcome, and the other three only under a unanimous split. The quantity that distinguishes the two readings, the
smallest discordant imbalance the test could resolve, is computable from the
discordant count alone and costs nothing to report. We suggest it accompany every
paired comparison on a small corpus.

Fourth, plan-level evaluation is a trade-off rather than an improvement, and one
term of that trade-off is a defect class rather than a cost. Because the compiled
plan is produced after expression resolution, policies observe concrete values
and avoid a class of interpolation-related imprecision that source-level analysis
must otherwise reason about \citep{opdebeeck2023control}. Plan generation in
exchange adds a step that the latency figures here exclude and that is expected
to dominate the layer's cost, and it requires a working provider configuration. It also introduces a distinction that a policy author must handle
explicitly, because Terraform serialises three configuration states of a scalar
optional attribute into shapes that a presence test orders wrongly. An attribute
set to a literal carries its value. An attribute set from a reference to another
resource is absent from the resolved attribute set and recorded as not-yet-known,
since the value is not resolvable until apply. An attribute not set at all is
present with the value null. A test for the attribute's presence therefore passes
the state that sets nothing and flags the state that sets it by reference, which
does not weaken such a control but inverts it: the compliant case is reported and
the violating case is let through. The two errors conceal each other in the
aggregate, since the resulting confusion matrix shows one false positive and one
false negative --- a signature indistinguishable from ordinary imprecision. We
encountered this in our own policy set and disclose the correction in
Section~\ref{sec:internal}. It is a property of plan-level evaluation rather than
of that implementation: any engine reading plan documents inherits the three-way
distinction, and any that collapses it inverts the same class of control. The
correct treatment is to report only the not-set state and to yield no finding on
the not-yet-known one, which is genuinely indeterminate at plan time. Block-typed
attributes are exempt, since an unconfigured block is likewise absent but a
configured one appears with its contents, so presence discriminates correctly
over blocks.

A fifth observation concerns tool succession rather than tool comparison. The
tfsec/Trivy pair measured here (Section~\ref{sec:rq2}) is a rule set carried from
a retired product into its maintained successor across a change of engine, and on
this corpus the two differ by four rules, all present in the predecessor and
silent in the successor; one of them costs the successor a detection. The magnitude is small; the direction is
what matters, since coverage was lost rather than gained and nothing in either
product's output announces it. Practitioners migrating between a scanner and its
replacement should not assume that coverage transfers rule for rule, and a
benchmark measuring only one of the two has no way to see the difference at all.

%------------------------------------------------------------------------------
\section{Threats to Validity}
\label{sec:threats}

Threats are organised under the categories of
\citet{wohlin2012experimentation}.

\subsection{Construct Validity}
\label{sec:construct}
The controlled corpus was authored by the same investigator who implemented the
evaluated policy set. Ground-truth labels therefore encode the same reading of
each control that the policies encode, and agreement between the two is not
independent evidence of detection capability. AI coding assistance was used in
writing both the policy set and the benchmark cases, so the two share that
influence as well as an author, and the dependence is if anything stronger than
single authorship alone implies. This is the most serious limitation
of the evaluation, and it is only partly mitigated.

Two mitigations apply. Labels are intended to derive from the text of the
originating control rather than from any policy implementation, though
Section~\ref{sec:labelling} documents one case where a label encoded the reference
policy's stricter reading instead, and both the specifications and the
generator that consumes them are released, so a reader can audit the derivation.
And the harness is also run on 25 third-party repositories the investigator did
not author, though these are unlabelled and support no effectiveness estimate
(Section~\ref{sec:external}).

A third check bears on it partially. The blind relabelling pass of
Section~\ref{sec:labelling} agreed with 47 of the 48 labels it covered, and the one
disagreement it surfaced was a genuine specification defect of exactly the
predicted kind: a control whose stated text was weaker than the requirement its
label encoded, where the stricter reading matched the reference policy set. That
the pass found such a case is evidence the mechanism is real rather than merely
conceivable.

The audit is not a sufficient mitigation. Its second labeller is a language model
whose reading of control texts derives from the same public documentation the
specifications were written from, so the two readings share provenance, agreement
is biased upward, and no claim of human inter-rater reliability is available from
it. It also ran in the same assistant session used for earlier debugging
(Section~\ref{sec:labelling}), so it is not independent of the AI assistance used
in building the study (Section~\ref{sec:environment}), which weakens its
independence further. Its
identity, parameters and instruction were not recorded, so the audit cannot be
repeated as run. A single investigator, with AI assistance, still wrote every
specification, and no labelled part of the corpus has a recorded source outside
this failure mode.

What would remove this threat rather than bound it is a separation of benchmark
authorship from policy implementation, and we state it as a protocol because this
corpus does not satisfy it: controls selected by someone other than the author of
the detection rules; cases generated without reference to any detection
implementation; labels fixed and recorded before the rules that detect them are
written; and the corpus frozen, by commit, before any tool is executed against it.
None holds fully here. The run manifest records the commit checked out at
measurement time, but the reported run was taken on a working tree with
uncommitted changes: the recorded commit is release 1.3.1, whose corpus had 44
internal cases, not the 52 measured. The measured corpus is therefore identified
by the released package (release 3.0.0) rather than by that commit, and the harness now records
the tree's uncommitted paths alongside it. The first three are the substantive
ones, and independent relabelling by a second human who has not read
the policy set is the cheapest step toward them. A benchmark satisfying all four
could offer its ground truth as independent evidence of detection capability;
this one offers it as a stated and auditable assumption.

Results on the controlled corpus should therefore not be read as an estimate of
field performance. A high score on a developer-authored corpus is at least as
consistent with insufficient corpus difficulty as with strong detection, and we
do not interpret it as the latter.

\subsection{Internal Validity}
\label{sec:internal}
Comparator tools ran with default rule sets, whereas the plan-level policies were
authored with knowledge of the corpus taxonomy. This asymmetry favours the
plan-level layer. It is disclosed rather than adjusted for, because equivalent
tuning effort across heterogeneous engines cannot be operationalised objectively;
readers should treat the comparison as bounding what default-configuration
scanning achieves, not as a ranking of the tools' attainable capability.

One plan-level policy was corrected after its results had been inspected, and the
correction is disclosed because it is exactly the iteration the comparators did
not receive. The audit-log-encryption rule tested for the presence of a
key-identifier attribute in the plan document, which inverts the control for the
reason set out in Section~\ref{sec:discussion}: the compliant case was reported as
violating and the violating case as compliant, and the two errors concealed each
other in the aggregate. The corrected rule distinguishes not-set from
not-yet-known and reports only the former.

Because this correction was informed by inspecting corpus results, the plan-level
figures reflect one round of iteration that Checkov, tfsec and Trivy did not
receive. The pre-correction and post-correction results are both in the
replication package.

Latency was measured on a single host with recorded specifications, so absolute
figures do not transfer across environments. Latency additionally measures each
tool at its own interface: the plan-level figure covers policy evaluation over an
already-compiled plan document and excludes the \texttt{terraform init} and
\texttt{terraform plan} invocations that produce it, which are expected to
dominate that layer's true wall-clock cost in continuous integration. The source-level scanners have no
comparable preparation step. Plan-level and source-level latencies are therefore
not interchangeable, and the plan-level figure is a lower bound on deployed cost.

\subsection{External Validity}
The corpus targets Terraform against a single cloud provider, and the native
module-testing layer is exercised against one project's infrastructure. Neither
the reported effectiveness nor the suite-size measurement of
Section~\ref{sec:rq3} generalises to other IaC technologies, other providers, or
the multi-module configurations typical of large organisations. No claim is made
about organisational factors such as policy ownership, escalation workflow, or
developer friction, because this study measures tool output and does not observe
practitioners.

The tool set bounds generalisation further: three third-party scanners are
measured but only two independent rule sets are represented
(Section~\ref{sec:toolsel}), so the comparison speaks to those rule sets rather
than to source-level IaC scanning in general. Provider coverage bounds it again,
since security-policy adoption in Terraform differs materially across cloud
providers \citep{verdet2025adoption} and this corpus targets one.

\subsection{Conclusion Validity}
The corpus is deterministic, not adversarial, so reported effectiveness
characterises performance on a published taxonomy of known misconfiguration
classes; no claim is made regarding novel or deliberately obfuscated insecure
configurations, or an adversary who has read the policy set. On statistical
conclusion validity, the $d_{\min}$ column of Table~\ref{tab:allpairs} bounds
what the paired tests could establish: under the Holm correction three of the six
comparisons among the four source- and plan-level tools could not have reached
significance under any outcome, and the other three only under a unanimous split. We therefore draw no conclusion of
equivalence anywhere in this paper, and the effectiveness comparison among
scanners should be read as unresolved rather than as settled in either direction.

%------------------------------------------------------------------------------
\section{Replication Package}
\label{sec:replication}

The replication package contains the benchmark corpus with ground-truth
annotations and the generator specifications from which those annotations derive;
the canonical control taxonomy as auditable data; the per-case output of the blind
relabelling pass of Section~\ref{sec:labelling}, with its reasons and its stated
limitations; the execution harness, which records raw tool output, resolved
versions, and repeated latency samples; the normalization engine; the statistical
implementation, with a test suite pinning each method to independently derived
values; and the analysis stage that generates every results table in this paper. It also
retains the measured results from before the plan-level policy correction
disclosed in Section~\ref{sec:internal}, so that the disclosure can be checked
rather than taken on trust.

Three properties are worth stating explicitly. The harness declines to emit
results when no case is admissible, instead of producing an empty but plausible
table. The statistical implementation refuses to report a proportion whose
denominator is zero, returning a not-estimable marker instead of a value a reader
could mistake for a measurement. And both diagrams in this paper are generated
from the recorded results rather than drawn, with a continuous-integration check
that fails if a re-measurement moves a corpus size or tool version that a figure
asserts.

Fig.~\ref{fig:artifact} gives the package layout. Ground-truth labels, the
scenario corpus, and the policy sources are separate top-level artefacts, so a
reader can substitute an independently authored label set without modifying the
harness.

Table~\ref{tab:repro} maps each reported result to the stage that emits it and the
recorded artefact it is computed from. One command,
\path{experiments/run_baselines.sh}, executes the scanners, writes the raw output,
and regenerates every row; \path{paper/Makefile} then refuses to typeset if any
generated table is absent, so a manuscript built from this package cannot contain a
results table that no recorded measurement supports; figures quoted in the prose
are computed from the same recorded output but are not checked mechanically. The mapping is given because
describing a package as complete is not the same as telling a reader which file
answers which question.

\begin{table}[!t]
\caption{Each reported result, the stage that emits it, and the recorded artefact
it derives from.}
\label{tab:repro}
\centering
\footnotesize
\setlength{\tabcolsep}{4pt}
\begin{tabularx}{\linewidth}{l l >{\raggedright\arraybackslash}X}
\toprule
\textbf{Result} & \textbf{Emitted by} & \textbf{Computed from} \\
\midrule
Corpus admissibility (Sec.~\ref{sec:admissibility}) & \texttt{corpus.py} &
\path{benchmark/}, \path{results/corpus_report.json} \\
\addlinespace[2pt]
Criterion sensitivity (RQ1) & \texttt{analyze.py} & \path{results/raw/} \\
\addlinespace[2pt]
Effectiveness and intervals (RQ2) & \texttt{analyze.py} & \path{results/raw/} \\
\addlinespace[2pt]
Paired tests and $d_{\min}$ (RQ2) & \texttt{analyze.py} &
\path{results/evaluation.json} \\
\addlinespace[2pt]
Latency and suite size (RQ3) & \texttt{analyze.py} &
\path{results/run_manifest.json}, \path{infrastructure/tests/} \\
\addlinespace[2pt]
Layer attribution (RQ4) & \texttt{analyze.py} & \path{results/raw/} \\
\addlinespace[2pt]
Labelled external cases (RQ5) & \texttt{analyze.py} & \path{results/raw/} \\
\addlinespace[2pt]
Unlabelled subset (RQ5) & \texttt{external.py} &
\path{results/raw/external/}, \path{results/external_subset.json} \\
\addlinespace[2pt]
Both diagrams & \texttt{generate\_figures.py} &
\path{results/evaluation.json}, \path{results/run_manifest.json} \\
\bottomrule
\end{tabularx}
\end{table}

% Set as a listing rather than a drawing. The diagram this replaces was a 10:1
% horizontal strip: at the column measure it rendered 0.34in tall with unreadable
% labels. A directory layout is conventionally a listing in any case, and
% verbatim carries it at body-text size with no scaling risk.
\begin{figure}[!t]
\centering
\footnotesize
\begin{verbatim}
benchmark/
  benchmark.json          planned case catalogue
  generate_corpus.py      case generator
  internal/               authored cases
  external/               CIS-derived cases,
                          unlabelled subset manifest
  labelling/              blind second pass
evaluation/
  corpus.py               admissibility gate
  normalize.py            finding normalization
  control_map.json        canonical taxonomy
  stats.py                exact CI + McNemar
  tables.py               LaTeX table emitters
  analyze.py              analysis entry point
  external.py             unlabelled subset
security_framework/
  policies/*.rego         plan-level policy set
  engine/                 layer 1 scanning
infrastructure/
  tests/*.tftest.hcl      native module suite
results/
  raw/                    per-case tool output
  run_manifest.json       resolved versions
  tables/                 generated tables
  pre_opa_polarity_fix/   pre-correction run
experiments/
  run_baselines.sh        full re-measurement
  generate_figures.py     figure sources
  fetch_external.py       pinned external repos
\end{verbatim}
% The caption said what the paragraph above it already says. Springer sets long
% captions in the same small type as the listing, so the duplicate was the least
% readable copy of the sentence; only the non-obvious directory is kept here.
\caption{Replication package layout. \texttt{pre\_opa\_polarity\_fix} retains the
measurements taken before the policy correction disclosed in
Section~\ref{sec:internal}.}
\label{fig:artifact}
\end{figure}

%------------------------------------------------------------------------------
\section{Conclusion}
\label{sec:conclusion}

This paper contributes an evaluation methodology for IaC security validation
tools: a finding-normalization scheme that makes heterogeneous scanner output
commensurable, an explicit treatment of matching strictness, a stated and
re-runnable corpus admissibility rule whose value as a filter is untested here,
and a statistical protocol suited to
paired small-sample comparison. The accompanying framework composes
repository-edge validation, native module testing, and plan-level policy
evaluation, and is released with a harness in which every reported figure is a
generated artefact traceable to recorded tool output.

The methodological findings are the transferable result. On this corpus the
matching criterion moves reported recall by 29 to 43 percent of the spread between
the tools compared, and by most of the available range for a tool whose findings
are of a different kind from the corpus's controls, which means comparative
claims in this area should report their criterion as part of their method.
Taxonomy coverage can favour the tool the taxonomy was authored around. Plan-level
policy evaluation inherits a three-way serialisation distinction that silently
inverts controls when it is collapsed. And corpus size is a weaker indicator of
evidential strength than interval width and the smallest difference the chosen
test could have resolved.

The outcome stated in Section~\ref{sec:intro} holds without adjustment: a
third-party scanner attains the highest point-estimate control-level recall and
the pipeline's scored layers do not lead the comparison. That is one indication
the procedure is not self-flattering, not proof of its neutrality, since
Section~\ref{sec:internal} discloses a tuning asymmetry running the other way.

Three lines of future work follow directly. A mapping between the security-smell
taxonomy that existing manually annotated IaC oracles are labelled against and
the CIS control identifiers this evaluation uses would allow effectiveness to be
measured on real-world configurations the investigator did not author, which is
one of the mitigations Section~\ref{sec:construct} identifies as outstanding and the
single change that would most strengthen the evidence. Enlarging the compliant
baseline would narrow the specificity intervals to a width at which
false-positive behaviour can be compared rather than merely reported. And
extending the taxonomy beyond Terraform would establish whether the criterion
sensitivity measured here is a property of IaC scanning or of this technology.

%------------------------------------------------------------------------------
% Springer's own heading for this block is "Statements and Declarations".
\section*{Statements and Declarations}

\paragraph{Funding} This research received no external funding.

\paragraph{Competing interests} The author designed and implemented the plan-level
policy set and the repository-edge layer evaluated here alongside third-party
tools. This is a non-financial interest, and Sections~\ref{sec:construct}
and~\ref{sec:internal} disclose the asymmetries it creates. The author declares no financial interests and is not
affiliated with, and receives no support from, the vendor or
maintainer of any tool evaluated in this study.

\paragraph{Ethics approval} Not applicable. This study analyses software
artefacts and involves no human participants or personal data.

\paragraph{Consent to participate} Not applicable.

\paragraph{Consent for publication} Not applicable.

\paragraph{Data and code availability} The corpus, ground-truth annotations,
canonical control taxonomy, normalization engine, execution harness, statistical
implementation, analysis code, and the raw per-case tool output from which every
table in this paper is generated are openly available. The version reported here
is release 3.0.0, archived at the first DOI below; the second always resolves to
the latest release.

\begin{center}
\url{https://doi.org/10.5281/zenodo.22921654}\\
\url{https://doi.org/10.5281/zenodo.21645016}\\
\url{https://github.com/mchittineni/IaCSecBench}
\end{center}

\paragraph{Author contributions} The sole author designed the study, implemented
the harness and the normalization engine, conducted the measurements, performed
the analysis, and wrote the manuscript, with AI assistance as described under Use
of AI-assisted technologies.

\paragraph{Use of AI-assisted technologies} AI assistance was used in three
roles. As a measuring instrument, a large language model was the second labeller
in the automated consistency audit of Section~\ref{sec:labelling}, which states
that the model's identity, parameters and instruction were not recorded and what
the resulting agreement does and does not establish in consequence; the per-case
labels and reasons are in the replication package. In building the study, AI
coding assistants were used to write parts of the harness and analysis code, the
plan-level policy set and the benchmark cases, and to help debug the harness and
inspect results, as described in Section~\ref{sec:environment}. In preparing the
manuscript, an AI assistant was used for language editing and to draft revisions
of manuscript passages under the author's direction. The assistant products used were Claude Code (Anthropic) and Gemini
(Google); their model versions and parameters were not recorded. The study design and conclusions are
the author's. The author reviewed all AI-assisted work and accepts full
responsibility for the manuscript and its content.

%------------------------------------------------------------------------------
\appendix
\section{Scope of the Validated Threats}
\label{app:threats}

The controls the pipeline enforces were selected against a STRIDE decomposition
(spoofing, tampering, repudiation, information disclosure, denial of service,
elevation of privilege) \citep{shostack2014threat} of a continuous-integration deployment path,
summarised in Table~\ref{tab:stride}. Protected assets are infrastructure
configurations and compiled plan states, ingested datasets, policy and test
definitions, and pipeline credentials; the trust boundaries separate the developer
workspace from version control, the evaluation runtime from the target
repository, and the compiled plan from the policy engine.

We record the decomposition to delimit which misconfiguration classes the corpus
exercises, and claim nothing further from it. The pipeline was not subjected to
adversarial testing and no adversary model was validated, so the table states
intended coverage rather than a result; it is not a threat model of IaC
pipelines, which this paper does not attempt.

\begin{table}[!t]
\caption{STRIDE categories, mitigations, and the layer that enforces them.}
\label{tab:stride}
\centering
\small
\setlength{\tabcolsep}{4pt}
% The mitigation column wraps, and a justified narrow column stretches its
% interword spaces to fill each line, which TeX reports as badly underfull.
% >{\raggedright\arraybackslash} sets it ragged-right instead; arraybackslash
% restores \\ as the row terminator, which \raggedright would otherwise rebind.
\begin{tabularx}{\linewidth}{l >{\raggedright\arraybackslash}X c}
\toprule
\textbf{Category} & \textbf{Mitigation} & \textbf{Layer} \\
\midrule
Spoofing & Authenticated API authorisers; transport encryption enforcement;
prohibition of embedded credentials & 1, 3 \\
\addlinespace[2pt]
Tampering & Schema validation; pre-commit enforcement; evaluation of compiled
plan documents & 1, 2 \\
\addlinespace[2pt]
Repudiation & Structured pipeline telemetry; recorded test execution; compliance
tagging & 2, 3 \\
\addlinespace[2pt]
Information disclosure & Pattern-based personal-data detection; mandatory
encryption at rest; storage public-access blocking & 1, 3 \\
\addlinespace[2pt]
Denial of service & Validation of web application firewall association, header
handling, and request throttling & 2, 3 \\
\addlinespace[2pt]
Elevation of privilege & Policy denial of wildcard identity actions and
unrestricted ingress ranges & 3 \\
\bottomrule
\end{tabularx}
\end{table}

Four domains are out of scope: post-deployment runtime compromise, including
container and hypervisor vulnerabilities; supply-chain compromise of provider
binaries or registry mirrors; IaC technologies other than Terraform, since the
plan-level layer presupposes Terraform plan JSON; and physical and
social-engineering attacks.

%------------------------------------------------------------------------------
% spbasic is Springer's author-year style and ships with the svjour3 template.
% plainnat is the closest author-year style available in a stock TeX
% distribution, so the fallback build produces the same citation form.
\IfFileExists{spbasic.bst}{\bibliographystyle{spbasic}}{\bibliographystyle{plainnat}}
\bibliography{refs}

\end{document}
